\begin{abstract}
\bi{Data agents repeatedly discover database-specific calculations and assumptions. Reusing this knowledge requires preserving the computation's meaning and maintaining its validity as data change. We present \system, a shared memory system that couples executable metric and feature definitions with data conditions and learning evidence. Deterministic publication checks validate extracted definitions before sharing. Compilation instantiates definitions for new questions and prediction instances; snapshot-bound metric execution checks their conditions on the consuming snapshot. Condition identities and dependency versions let agents share validation, repair, and optimization work. On a TPC-DS workload with \PpAgents\ agents, \system\ reduces tokens by \PpMavraTokSaveNone\% relative to no memory and by \PpMavraTokSavePrivate\% relative to private memory with the same mechanisms. After a data restatement, accuracy is \PpMavraChangeAcc\% versus \PpNoneChangeAcc\% without memory. In a separate official-API RelBench study, total tokens relative to no memory decrease by \OfStackSaveAll\% on rel-stack and \OfTrialSaveAll\% on rel-trial; during reuse, the decreases are \OfStackSaveReuse\% and \OfTrialSaveReuse\%. Five of six tasks meet predeclared full-sequence quality margins; one regression task has \OfAdverseRegression\% higher MAE. The results quantify reduced repeated exploration and shared maintenance, together with their quality and validity limits.}
{数据智能体反复发现数据库特有的计算逻辑与假设。复用这些知识需要保留计算含义，并在数据变化时维持有效性。本文提出共享记忆系统\system，将可执行的指标与特征定义同数据条件和学习证据结合。确定性的发布检查在共享前验证提取结果；编译为新问题与预测样本实例化定义，快照绑定的指标执行在消费快照上检验条件。条件标识与依赖版本使智能体共享验证、修复与优化工作。在\PpAgents{}个智能体的TPC-DS负载上，\system\ 相对无记忆减少\PpMavraTokSaveNone\%的token，相对采用相同机制的私有记忆减少\PpMavraTokSavePrivate\%。数据更正后，正确率为\PpMavraChangeAcc\%，无记忆为\PpNoneChangeAcc\%。在独立的官方API RelBench实验中，相对无记忆，rel-stack与rel-trial的总token分别减少\OfStackSaveAll\%与\OfTrialSaveAll\%，复用阶段分别减少\OfStackSaveReuse\%与\OfTrialSaveReuse\%。六项任务中五项满足预设的全序列质量门槛，一项回归任务的MAE高\OfAdverseRegression\%。结果量化了重复探索与共享维护的节省，同时给出质量与有效性边界。}
\end{abstract}

\begin{IEEEkeywords}
\bt{data agents, shared memory, metric definitions, feature engineering, data validation}
{数据智能体，共享记忆，指标定义，特征工程，数据验证}
\end{IEEEkeywords}
